\documentclass[12 pt, letterpaper]{hmcpset}
\usepackage{hw-preamble}

\name{Jayden Thadani}
\class{Principles of Mathematical Analysis self-study}
\assignment{Chapter 1 exercises --- The Real and Complex Number Systems}
\duedate{}

\begin{document}

\begin{problem}[1.]
	If $r$ is rational ($r \neq 0$) and $x$ is irrational prove that $r + x$ and $rx$ are irrational.
\end{problem}

\begin{solution} ~
	
	Suppose, by way of contradiction, that we had $0 \neq r \in \mathbb{Q}$ and $r + x \in \mathbb{Q}$ or $rx \in \mathbb{Q}$ for some $x \in \mathbb{R} \setminus \mathbb{Q}$. Then we would have to have $x = r + x - x \in \mathbb{Q}$ or $x = rx/r \in \mathbb{Q}$. Each of these directly contradicts $x \notin \mathbb{Q}$. Thus, $r + x$ and $rx$ must both be irrational.
\end{solution}

\pagebreak

\begin{problem}[2.]
	Prove that there is no rational number whose square is $12$.
\end{problem}

\begin{solution} ~
	
	Suppose, by way of contradiction, that there is some rational number $p$ with $p^{2} = 12$. Then there must be some $m, n \in \mathbb{Z}$ such that $p = m/n$ and $\gcd (m, n) = 1$. Thus, we must have
	\[
		\begin{split}
			m^{2} & = 12 n^2 \\
			      & = 3 (2n)^{2}.
		\end{split}
	\]
	If $3 \mid m^{2}$, we must have $3 \mid m$ since $3$ is prime. But then, we can write $m = 3r$ for some $r \in \mathbb{Z}$ giving us
	\[
		9 r^{2} = 3 (2n)^{2}
	\]
	and thus,
	\[
		4n^{2} = 3 r^{2}.
	\]
	Since $3 \nmid 4$, we must have $3 \mid n^{2}$ and thus, $3 \mid n$. This contradicts the fact that $\gcd(m, n) = 1$, and thus, there cannot exist any rational number $p$ with $p^{2} = 12$.
\end{solution}

\pagebreak

\begin{problem}[2*.]
	Prove that between any two real numbers there is an irrational number.
\end{problem}

\begin{solution}~
	
	We know (by Theorem 1.20) that there is a rational number between any two reals. Thus, for any two $a, b \in \mathbb{R}$ with $a < b$, and a positive $x \in \mathbb{R} \setminus \mathbb{Q}$, we have some rational $r$ such that
	\[
		a/x < r < b/x.
	\]
	Noting that $rx$ must be irrational (shown in exercise 1), we have
	\[
		a < rx < b
	\]
	giving us the irrational number between $a$ and $b$.
\end{solution}

\pagebreak

\begin{problem}[3.]
	Prove Proposition 1.15 --- that the axioms for multiplication imply the following statements
	\begin{enumerate}
		\item If $x \neq 0$ and $xy = xz$ then $y = z$.
		\item If $x \neq 0$ and $xy = x$ then $y = 1$.
		\item If $x \neq 0$ and $xy = 1$, then $y = 1/x$.
		\item If $x \neq 0$, then $1/(1/x) = x$.
	\end{enumerate}
\end{problem}

\begin{solution}
	\begin{enumerate}
		\item Since $x \neq 0$, we can multiply both sides by $1/x$ to get the following series of equivalent equalities.
			\[
				\begin{split}
					\frac{1}{x} xy & = \frac{1}{x} xz. \\
					1y & = 1z. \\
					y & = z.
				\end{split}
			\]
		\item Again, we can multiply by $1/x$ to get equivalent equalities ---
			\[
				\begin{split}
					xy & = x. \\
					\frac{1}{x} xy & = \frac{1}{x} x. \\
					1 y & = 1. \\
					y & = 1.
				\end{split}
			\]
		\item Multiply by $1/x$ ---
			\[
				\begin{split}
					\frac{1}{x} xy & = \frac{1}{x} 1. \\
					1 y & = \frac{1}{x}. \\
					y & = \frac{1}{x}.
				\end{split}
			\]
		\item Since $(1/x) x = 1$, $x$ is the unique multiplicative inverse of the $(1/x)$. Thus, $x = (1/(1/x))$.
	\end{enumerate}
\end{solution}

\pagebreak

\begin{problem}[4.]
	Let $E$ be a nonempty subset of an ordered set; suppose $\alpha$ is a lower bound of $E$ and $\beta$ is an upper bound of $E$. Prove that $\alpha \le \beta$.
\end{problem}

\begin{solution}
	Choose any $x \in E$. Then $\alpha \le x \le \beta$, and thus, by transitivity, $\alpha \le \beta$.
\end{solution}

\pagebreak

\begin{problem}[5.]
	Let $A$ be a nonempty set of real numbers which is bounded below. Let $-A$ be the set of all numbers $-x$ where $x \in A$. Prove that
	\[
		\inf A = - \sup (-A).
	\]
\end{problem}

\begin{solution}~
	
	Note that any lower bound $\alpha$ of $A$ gives $-\alpha \ge -a$ (by reversing signs from $\alpha \le a$). Thus, any lower bound of $A$ is an upper bound of $-A$. \\
	
	Then, by definition, for any $a \in A$, and any lower bound $\alpha$ of $A$, we have
	 \[
		-a \le \sup(-A) \le -\alpha.
	\]
	Negating and reversing the signs gives us
	\[
		\alpha \le -\sup(-A) \le a.
	\]
	That is, by definition, $-\sup(-A)$ is $\inf A$.
\end{solution}

\pagebreak

\begin{problem}[6.]
	Fix $b > 1$.
	\begin{enumerate}
		\item If $m, n, p, q$ are integers,  $n > 0$, $q > 0$ and $r = m/n = p/q$, prove that
			\[
				(b^{m})^{1/n} = (b^{p})^{1/q}.
			\]
			Hence it makes sense to define $b^{r} = (b^{m})^{1/n}$.
		\item Prove that $b^{r + s} = b^{r}b^{s}$ if $r$ and $s$ are rational.
		\item If $x$ is real, define $B(x)$ to be the set of all numbers $b^{t}$ where $t$ is rational and $t \le x$. Prove that
			\[
				b^{r} = \sup B(r)
			\]
			when $r$ is rational. Hence it makes sense to define
			\[
				b^{x} = \sup B(x)
			\]
			for every real $x$.
		\item Prove that $b^{x + y} = b^{x} b^{y}$ for all real $x$ and $y$.
	\end{enumerate}
\end{problem}

\begin{solution}
	\begin{enumerate}
		\item If $m/n = p/q$, we have $mq = pn = k$. Thus, we can write $(b^{m})^{1/n}$ and $(b^{p})^{1/q}$ as $k$th roots of $b^{mp}$ ---
			\[
				\begin{split}
					(b^{mp})^{1/k} & = (((b^{m})^{p})^{1/p})^{1/n} \\
						       & = (b^{m})^{1/n}
				\end{split}
			\]
			and
			\[
				\begin{split}
					(b^{mp})^{1/k} & = (((b^{p})^{m})^{1/m})^{1/q} \\
						       & = (b^{p})^{1/q}.
				\end{split}
			\]
			Thus, by the uniqueness of $k$th roots, we have $(b^{m})^{1/n} = (b^{p})^{1/q}$. \\

			We used the following result in the proof above.

			\begin{lemma}
				For any $1 < b \in \mathbb{R}$ and $x, y \in \mathbb{N}$, $b^{1/xy} = (b^{1/x})^{1/y}$.
			\end{lemma}

			\begin{proof}
				Notice that for any $a \in \mathbb{R}$, $a^{xy} = (a^{y})^{x}$. Thus,
				\[
					\begin{split}
						((a^{xy})^{1/x})^{1/y} & = (((a^{y})^{x})^{1/x})^{1/y} \\
								       & = (a^{y})^{1/y} \\
								       & = a.
					\end{split}
				\]
				For any $b \in \mathbb{R}$, choose $a = b^{1/xy}$. Thus, $b = a^{xy}$. Then $b^{1/xy} = a$ and $(b^{1/x})^{1/y} = a$. That is, $b^{1/xy} = (b^{1/x})^{1/y}$.
			\end{proof}
		\item Let $r = r_{1}/r_{2}$ and $s = s_{1}/s_{2}$. Then $r + s = (r_{1} s_{2} + s_{1} r_{2})/r_{2} s_{2}$. Then
			\[
				\begin{split}
					b^{r + s} & = b^{(r_{1} s_{2} + s_{1} r_{2})/r_{2} s_{2}} \\
						  & = (b^{r_{1} s_{2} + s_{1} r_{2}})^{1/r_{2} s_{2}} \\
						  & = (b^{r_{1} s_{2}} b^{s_{1} r_{2}})^{1/r_{2} s_{2}} \\
						  & = (b^{r_{1} s_{2}})^{1/r_{2} s_{2}} (b^{s_{1} r_{2}})^{1/r_{2} s_{2}} \\
						  & = b^{r_{1} s_{2}/r_{2} s_{2}} b^{s_{1} r_{2}/s_{2} r_{2}} \\
						  & = b^{r_{1}/r_{2}} b^{s_{1}/s_{2}} \\
						  & = b^{r} b^{s}.
				\end{split}
			\]
			as desired.
		\item First we will show that $b^{r}$ is an upper bound on $B(r)$. \\

			Any element of $B(r)$ is of the form $b^{t}$ for some rational $t \le r$. We can write $b^{r} = b^{t} b^{r - t}$ where $r - t$ is some nonnegative rational. Since $b > 1$, $b^{r - t} \ge 1$, and thus, $b^{r} \ge b^{t}$. Thus, $b^{r}$ is an upper bound on $B(r)$. \\
			
			Now we will show that it is the least upper bound. \\

			Suppose $b^{r} > \beta \in \mathbb{R}$. Then $\beta$ is less than an element of  $B(r)$ --- $b^{r}$ itself, and thus, is not an upper bound. That is, no real less than $b^{r}$ can be an upper bound of $B(r)$ --- $b^{r}$ is the least upper bound of $B(r)$.
		\item First we will show that $b^x b^{y}$ is an upper bound on $B(x + y)$. \\

			Note that any element of $B(x + y)$ can be written as $b^{q}$ for some $q \in \mathbb{Q}$ with $q \le x + y$. In turn, $q$ can give us two rationals $q_{1}$ and $q_{2}$, with $q_{1} \le x$ and $q_{2} \le y$ and $q_{1} + q_{2} \ge q$. Choose $q_{1}$ to be a rational between $x - (x + y - q)/2$ and $x$. Choose $q_{2}$ to be a rational between $y - (x + y - q)/2$ and $y$. Thus, $b^{q} \le b^{q_{1}} b^{q_{2}} \le b^{x} b^{y}$, giving us that $b^{x} b^{y}$ is an upper bound. \\

			Now we show that it is the least upper bound. \\

			Suppose $b^{x} b^{y} > \beta \in \mathbb{R}^{+}$. Then $\beta/b^{x} < r < b^{y}$ and where $r$ is a rational between the two reals. Since $\beta/r < b^{x}$, it is not an upper bound of $B(x)$. That is, there is some rational $t_{1} \le x$ with $\beta/r < b^{t_{1}}$. Similarly, since $r < y$, it is not an upper bound of $B(y)$, and thus, $r < b^{t_{2}}$ for some rational $t_{2} \le y$. Then $t_{1} + t_{2}$ is a rational less than $x + y$ with $\beta \le b^{t_{1} + t_{2}}$ ---
			\[
				\begin{split}
					\beta & = \frac{\beta}{r} r \\
					      & \le b^{t_{1}} b^{t_{2}} \\
					      & = b^{t_{1} + t_{2}}.
				\end{split}
			\]
	\end{enumerate}
\end{solution}

\pagebreak

\begin{problem}[7.]
	Fix $b > 1$, $y > 0$ and prove that there is a unique real $x$ such that $b^{x} = y$, by completing the following outline. (This $x$ is called the \textit{logarithm of $y$ to the base $b$}.)
	\begin{enumerate}
		\item For any positive integer $n$, $b^{n} - 1 > n(b - 1)$.
		\item Hence $b - 1 > n(b^{1/n} - 1)$.
		\item If $t > 1$ and $n > (b - 1)/(t - 1)$, then $b^{1/n} < t$.
		\item If $w$ is such that $b^{w} < y$, then $b^{w + 1/n} < y$ for sufficiently large $n$.
		\item If $b^{w} > y$, then $b^{w - 1/n} > y$ for sufficiently large $n$.
		\item Let $A$ be the set of all $w$ such that $b^{w} < y$ and show that $x = \sup A$ satisfies $b^{x} = y$.
		\item Prove that this $x$ is unique.
	\end{enumerate}
\end{problem}

\begin{solution}~
	\begin{enumerate}
		\item Note that $b - 1$ divides $b^{n} - 1$. Specifically,
			\[
				b^{n} - 1 = (b^{n - 1} + b^{n - 2} + \cdots + b + 1)(b - 1).
			\]
			But since, $b > 1$, each $b^{k} > 1$ (for $k \in \mathbb{N}_{n - 1}$). That gives us
			\[
				\begin{split}
					b^{n} - 1 & > (\underbrace{1 + \cdots + 1}_{n - 1 \text{ times}} + 1)(b - 1) \\
						  & = n(b - 1)
				\end{split}
			\]
			as desired.
		\item Let $a = b^{1/n}$. Note that $a > 1$; else we wouldn't have $a^{n} = b > 1$. Thus, we have $a^{n} - 1 > n(a - 1)$. This is just $b - 1 > n(b^{1/n} - 1)$.
		\item We already have $b - 1 > n(b^{1/n} - 1)$. Since $n > (b - 1)/(t - 1)$, we have
			\[
				b - 1 > \frac{b - 1}{t - 1}(b^{1/n} - 1).
			\]
			Rearranging, we get
			\[
				b^{1/n} - 1 < t - 1
			\]
			and finally $b^{1/n} < t$.
		\item Choose $t = y/b^{w} > 1$. Then by (c), for sufficiently large $n$ we can get $b^{1/n} < t$. Thus,
			\[
				\begin{split}
					b^{w + 1/n} & = b^{w} b^{1/n} \\
						    & < b^{w} t \\
						    & = b^{w} \frac{y}{b^{w}} \\
						    & = y.
				\end{split}
			\]
		\item Choose $t = b^{w}/y > 1$ and proceed similarly to (d).
		\item Suppose by way of contradiction that $x$ does not give $b^{x} = y$. We must have $b^{x} > y$ or $b^{x} < y$. We will show neither is possible. \\

			If we had $b^{x} > y$, we would have $b^{x - 1/n} > y$ for some large enough $n \in \mathbb{N}$. Then since $x - 1/n$ would be an upper bound of $A$ and $x - 1/n < x$, $x$ would not be the least upper bound of $A$. \\

			If we had $b^{x} < y$, we would have $b^{x + 1/n} < y$ for some large enough $n \in \mathbb{N}$. Then since $x + 1/n \in A$  and $x + 1/n > x$, $x$ would not be an upper bound of $A$.
		\item This follows from
			\begin{lemma}
				For any $1 < b \in \mathbb{R}$ and $x_{1} < x_{2} \in \mathbb{R}^{+}$, $b^{x_{1}} < b^{x_{2}}$.
			\end{lemma}
			\begin{proof}
				First we show that the result holds for rationals. Suppose we have $p_{1}, p_{2} \in \mathbb{Q}$ with $p_{1} < p_{2}$. Write $p_{1} = m_{1}/n$ and $p_{2} = m_{2}/n$ for $m_{1}, m_{2}, n \in \mathbb{Z}$ with $m_{1} < m_{2}$. Then $b^{m_{1}} < b^{m_{2}}$. But then, we must have $(b^{m_{1}})^{1/n} < (b^{m_{2}})^{1/n}$. Else, on raising to the power of $n$, we'd get $b^{m_{2}} \le b^{m_{1}}$. Thus, $b^{p_{1}} < b^{p_{2}}$.
				
				Now we can show the corresponding result for positive reals $x_{1} < x_{2}$. Choose rationals $q_{1}, q_{2}$ with $x_{1} < q_{1} < q_{2} < x_{2}$. Then $b^{q} < b^{q_{1}} < b^{q_{2}}$ for any rational $q$ with $q < x_{1} < q_{1} < q_{2}$. Thus, $b^{q_{1}}$ is an upper bound of $B(x_{1})$. Also, by definition, $b^{q_{2}} \in B(x_{2})$ and thus, $b^{q_{2}} \le b^{x_{2}}$. Then, $b^{x_{1}} \le b^{q_{1}} < b^{q_{2}} \le b^{x_{2}}$, which gives us $b^{x_{1}} < b^{x_{2}}$.
			\end{proof}
			Thus, if $x, x' \in \mathbb{R}^{+}$ give $b^{x} = b^{x'} = y$, we must have $x = x'$. If we didn't we'd have $b^{x} < b^{x'}$ or $b^{x'} < b^{x}$.
	\end{enumerate}
\end{solution}

\pagebreak

\begin{problem}[8.]
	Prove that no order can be defined in the complex field that turns it into an ordered field.
\end{problem}

\pagebreak

\begin{problem}[9.]
	Suppose $z = a + bi$, $w = c + di$. Define $z < w$ if $a < c$, and also if $a = c$ but $b < d$. Prove that this turns the set of all complex numbers into an ordered set. (This type of order relation is called a \textit{dictionary order} or \textit{lexicographic order}, for obvious reasons.) Does this ordered set have the least upper bound property?
\end{problem}

\pagebreak

\begin{problem}[10.]
	Suppose $z = a + bi$, $w = u + vi$, and
	\[
		a = \left ( \frac{\lvert w \rvert + u}{2} \right )^{1/2} \quad b = \left ( \frac{\lvert w \rvert - u}{2} \right )^{1/2}.
	\]
	Prove that $z^{2} = w$ if $v \ge 0$ and that $\overline{z}^{2} = w$ if $v \le 0$. Conclude that every complex number (with one exception) has two complex square roots.
\end{problem}

\pagebreak

\begin{problem}[11.]
	If $z$ is a complex number, prove that there exists an $r \ge 0$ and a complex number $w$ with $\lvert w = 1 \rvert$ such that $z = r w$. Are $w$ and $r$ always uniquely determined by $z$?
\end{problem}

\pagebreak

\begin{problem}[12.]
	If $z_{1}, \dots, z_{n}$ are complex, prove that
	\[
		\lvert z_{1} + z_{2} + \dots + z_{n} \rvert \le \lvert z_{1} \rvert + \lvert z_{2} \rvert + \dots + \lvert z_{n} \rvert.
	\]
\end{problem}

\pagebreak

\begin{problem}[13.]
	If $x, y$ are complex, prove that
	\[
		\lvert \lvert x \rvert - \lvert y \rvert \rvert \le \lvert x - y \rvert.
	\]
\end{problem}

\pagebreak

\begin{problem}[14.]
	If $z$ is a complex number such that $\lvert z  \rvert = 1$, compute
	\[
		\lvert 1 + z^{2} \rvert + \lvert 1 - z^{2} \rvert.
	\]
\end{problem}

\pagebreak

\begin{problem}[15.]
	Under what conditions does equality hold in the Schwarz inequality?
\end{problem}

\pagebreak

\begin{problem}[16.]
	Suppose $k \ge 3$, $\mathbf{x}, \mathbf{y} \in \mathbb{R}^{k}$, $\lVert \mathbf{x} - \mathbf{y} \rVert = d > 0$, and $r > 0$. Prove:
	\begin{enumerate}
		\item If $2r > d$, there are infinitely many $\mathbf{z} \in \mathbb{R}^{k}$ such that 
			\[
				\lVert \mathbf{z} - \mathbf{x} \rVert = \lvert \mathbf{z} - \mathbf{y} \rvert = r.
			\]
		\item If $2r = d$, there is exactly one such $\mathbf{z}$.
		\item If $2r < d$, there is no such $\mathbf{z}$.
	\end{enumerate}
	How must these statements be modified if $k$ is $2$ or $1$?
\end{problem}

\pagebreak

\begin{problem}[17.]
	Prove that
	\[
		\lVert \mathbf{x} + \mathbf{y} \rVert^{2} + \lVert \mathbf{x} - \mathbf{y} \rVert^{2} = 2 \lVert \mathbf{x} \rVert^{2} + 2 \lVert \mathbf{y} \rVert^{2}
	\]
	f $\mathbf{x} \in \mathbb{R}^{k}$ and $\mathbf{y} \in \mathbb{R}^{k}$. Interpret this geometrically, as a statement about parallelograms.
\end{problem}

\pagebreak

\begin{problem}[18.]
	If $k \ge 2$ and $\mathbf{x} \in \mathbb{R}^{k}$, prove that there exists $\mathbf{y} \in \mathbb{R}^{k}$ such that $\mathbf{y} \neq \mathbf{0}$, but $\mathbf{x} \cdot \mathbf{y} = 0$. Is this also true if $k = 1$?
\end{problem}

\pagebreak

\begin{problem}[19.]
	Suppose $\mathbf{a} \in \mathbb{R}^{k}$ and $\mathbf{b} \in \mathbb{R}^{k}$. Find $\mathbf{c} \in \mathbb{R}^{k}$ and $r > 0$ such that
	 \[
		\lVert \mathbf{x} - \mathbf{a} \rVert = 2 \lVert \mathbf{x} - \mathbf{b} \rVert
	\]
	if and only if $\lVert \mathbf{x} - \mathbf{c} \rVert = r$.
\end{problem}

\pagebreak

\begin{problem}[20.]
	With reference to the Appendix, suppose that property (III) were omitted from the definition of a cut. Keep the same definitions of order and addition. Show that the resulting ordered set has the least-upper-bound property, that addition satisfies axioms (A1) to (A4) (with a slightly different zero-element) but that (A5) fails.
\end{problem}

\end{document}
